% Bai 01
\chapter{Bài 1: Cài đặt và làm quen môi trường lập trình Python}

\section{Bài 1.1: Cài đặt Python dùng bản phân phối Anaconda}
\subsection{Mô tả bài tập và ý tưởng giải quyết}
\begin{itemize}
  \item Tải và cài đặt Anaconda/Miniconda.
  \item Kiểm tra phiên bản Python sau khi cài đặt.
  \item Ghi nhận môi trường đã thiết lập để phục vụ thực hành.
\end{itemize}

\subsection{Code cài đặt}
Phần này chủ yếu thao tác cài đặt (không có code Python). Nếu có lệnh kiểm tra, có thể ghi lại:

\begin{lstlisting}
python --version
python
\end{lstlisting}

\subsection{Kết quả chạy và giải thích}
Chèn ảnh chụp màn hình kết quả kiểm tra phiên bản/console.

\ScreenshotFigure{assets/screenshots/bai01_b1_python_version.png}{Kiểm tra phiên bản Python sau khi cài đặt (screenshot).}

\textbf{Giải thích:} Ảnh thể hiện Python đã được cài đặt và có thể chạy từ terminal/console.

\section{Bài 1.2: Quản trị môi trường Anaconda}
\subsection{Mô tả bài tập và ý tưởng giải quyết}
\begin{itemize}
  \item Liệt kê các môi trường conda, liệt kê packages.
  \item Cập nhật conda/pip và cài các thư viện cần thiết (ipykernel, jupyter/jupyterlab, matplotlib, scipy, scikit-learn).
\end{itemize}

\subsection{Code cài đặt}
\begin{lstlisting}
conda env list
conda list
python -m pip install --upgrade pip
conda update conda
conda install ipykernel jupyter
pip install jupyterlab==3.1.12
pip install matplotlib scipy scikit-learn
\end{lstlisting}

\subsection{Kết quả chạy và giải thích}
\ScreenshotFigure{assets/screenshots/bai01_b2_conda_env_list.png}{Liệt kê môi trường conda và packages (screenshot).}

\textbf{Giải thích:} Môi trường \texttt{base} (hoặc môi trường thực hành) hoạt động và các thư viện đã được cài đặt thành công.

\section{Bài 1.3: Lập trình trên JupyterLab/Jupyter Notebook}
\subsection{Mô tả bài tập và ý tưởng giải quyết}
\begin{itemize}
  \item Khởi động JupyterLab tại thư mục thực hành.
  \item Tạo notebook, viết markdown + code để in \texttt{Hello, Python World!}.
  \item Chạy thử bằng Python console / notebook và ghi lại kết quả.
\end{itemize}

\subsection{Code cài đặt}
\begin{lstlisting}
print("Hello, Python World!")
\end{lstlisting}

\subsection{Kết quả chạy và giải thích}
\ScreenshotFigure{assets/screenshots/bai01_b3_hello_world_notebook.png}{Notebook in ra \texttt{Hello, Python World!} (screenshot).}

\textbf{Giải thích:} Lệnh \texttt{print} xuất chuỗi ra màn hình, xác nhận môi trường notebook chạy được code Python.

\clearpage

